%%%%%%%%%%%%%%%%%%%%%%%%%%%%%%%%%%%%%%%%%
% Beamer Presentation - LaTeX Template
% Version 2.0 (March 8, 2022)
% Original Template: https://www.LaTeXTemplates.com
% Author: Vel (vel@latextemplates.com)
% License: CC BY-NC-SA 4.0

% Este modelo de apresentação foi 
% criado a partir do modelo de Giovanni Spadaro.
% Disponível em: https://github.com/Giovo17/presentation-template-unict-lm-data
%
% Adaptado por Lucas Amaral Taylor para criar uma versão especial 
% para os alunos de Matemática e Estatística da USP (IME-USP).
% Disponível em: https://github.com/lucasamtaylor01/IME-template
%%%%%%%%%%%%%%%%%%%%%%%%%%%%%%%%%%%%%%%%%

%----------------------------------------------------------------------------------------
% CLASSE DO DOCUMENTO E CONFIGURAÇÕES BÁSICAS
%----------------------------------------------------------------------------------------
\documentclass[
    11pt,               % Tamanho padrão da fonte
    % t,                % Alinhar verticalmente ao topo
    %aspectratio=169,   % Definir proporção 16:9
]{beamer}
\graphicspath{{img/}}         % Define o diretório das imagens

%----------------------------------------------------------------------------------------
% PACOTES NECESSÁRIOS
%----------------------------------------------------------------------------------------
\usepackage{
    booktabs,     % Melhora a aparência das linhas em tabelas
    palatino,     % Define Palatino como fonte principal
    subcaption    % Suporte para subfiguras
}
\usepackage[default]{opensans}  % Define Open Sans como fonte secundária
\input{config/code_langs}       % Importa configurações para highlight de código

%----------------------------------------------------------------------------------------
% CONFIGURAÇÃO DO TEMA
%----------------------------------------------------------------------------------------
% Tema Base
\usetheme{Boadilla}                          % Define o tema principal
\useinnertheme{circles}                      % Tema interno com círculos
\useoutertheme{miniframes}                   % Tema externo com miniframes
\setbeamertemplate{navigation symbols}{}     % Remove símbolos de navegação

% Cores Personalizadas
\definecolor{primaryColor}{RGB}{40,96,128}   % Cor primária - azul escuro
\definecolor{secondaryColor}{RGB}{217,217,217} % Cor secundária - azul médio
\definecolor{thirdColor}{RGB}{18,80,105} % Cor secundária - azul médio

% Configurações de Cores
\setbeamercolor{structure}{fg=primaryColor}
\setbeamercolor{palette primary}{bg=primaryColor, fg=white}
\setbeamercolor{palette secondary}{bg=secondaryColor, fg=white}
\setbeamercolor{title}{bg=primaryColor, fg=white}

% Cores do Cabeçalho e Rodapé
\setbeamercolor{headline}{bg=secondaryColor, fg=white}
\setbeamercolor{section in head/foot}{bg=primaryColor, fg=white}
\setbeamercolor{subsection in head/foot}{bg=secondaryColor, fg=white}
\setbeamercolor{author in head/foot}{bg=primaryColor, fg=white}
\setbeamercolor{title in head/foot}{bg=thirdColor, fg=white}
\setbeamercolor{date in head/foot}{bg=primaryColor, fg=white}
\setbeamercolor{page number in head/foot}{bg=thirdColor, fg=white}

%----------------------------------------------------------------------------------------
% BIBLIOGRAFIA
%----------------------------------------------------------------------------------------
% \usepackage[style=alphabetic,backend=biber]{biblatex}
% \addbibresource{bibliografia.bib}

\usepackage[backend=biber,style=numeric]{biblatex}
\addbibresource{bibliografia.bib}

%----------------------------------------------------------------------------------------
% HIPERLINKS
%----------------------------------------------------------------------------------------
\usepackage{hyperref}
\hypersetup{
  hidelinks,
  urlcolor=blue
}
%----------------------------------------------------------------------------------------
% INFORMAÇÕES DA APRESENTAÇÃO
%----------------------------------------------------------------------------------------
\title[Report]{Report}          % [Versão curta]{Versão completa}
\author[Nara Avila]{Nara Avila Moraes}            % [Versão curta]{Nome completo}
\institute[IFUSP]{Physics Institute \\ IFUSP}
\date[2026]{2026 / 01 / 16}

%----------------------------------------------------------------------------------------
% INÍCIO DO DOCUMENTO
%----------------------------------------------------------------------------------------
\begin{document}

% Slide de título com logo
\begin{frame}
    \begin{figure}
        \includegraphics[width=0.45\linewidth]{img/logo_IFUSP.png}
    \end{figure}
    \titlepage
\end{frame}

% Sumário
\begin{frame}
    \frametitle{Summary}
    \tableofcontents
\end{frame}

% Inclusão das seções
\section{Tasks Status} % Seções são adicionadas para organizar sua apresentação em blocos discretos, todas as seções e subseções são automaticamente exibidas no índice como uma visão geral da apresentação, mas NÃO são exibidas como slides separados.

%----------------------------------------------------------------------------------------

\begin{frame}
  \frametitle{Tasks Status}
    \begin{table}[h]
        \centering
        \begin{tabular}{|l|c|c|}
        \hline
        \textbf{Task} & \textbf{Status} & \textbf{Link} \\
        \hline
         &  & \\
        \hline
         &  & \\
        \hline
        \end{tabular}
    \end{table}
\end{frame}

%----------------------------------------------------------------------------------------
% \begin{frame}{Uma imagem}
%     \frametitle{Historical Records}
%     In addition, it is possible to retrieve the full search history by running the script directly via GitHub Actions.
%    \begin{figure}[h]
%        \centering
%        \includegraphics[width=0.3\textwidth]{img/github_topic_history.png}
%        \caption{GitHub Actions interface illustrating how the workflow can be executed to retrieve the full search history for a specific topic.}
%        \label{fig:github_topic_history}
%    \end{figure}
% \end{frame}
%----------------------------------------------------------------------------------------

% \begin{frame}
% 	\frametitle{Texto em tópicos}
%      Lorem ipsum dolor sit amet, consectetur adipiscing elit:
%     \begin{itemize}
%         \item Lorem ipsum dolor sit amet.
%         \item Lorem ipsum dolor sit amet.
%     \end{itemize}
	
% \end{frame}



\section{Overview of Learned Concepts} % Seções são adicionadas para organizar sua apresentação em blocos discretos, todas as seções e subseções são automaticamente exibidas no índice como uma visão geral da apresentação, mas NÃO são exibidas como slides separados.

%----------------------------------------------------------------------------------------

\begin{frame}
  \frametitle{ Recent studies}
  \begin{table}
    \begin{tabular}{|p{0.58\textwidth}|c|p{0.22\textwidth}|}
      \hline
      \textbf{Article} & \textbf{Status} & \textbf{Link} \\
      \hline
      Characterizing the non-equilibrium quark-gluon plasma with photons and hadrons & \inprogress & \href{https://j-f-paquet.github.io/theses/phd_thesis_final_JFPaquet.pdf}{Thesis} \\
      \hline
      Evidence for a New State of Matter   & \done & \href{https://arxiv.org/pdf/nucl-th/0002042}{Link} \\
      \hline
      Bayesian parameter estimation for  relativistic heavy-ion collisions & \inprogress & \href{https://arxiv.org/pdf/1804.06469}{Thesis} \\
      \hline
       Small System Collectivity  in Relativistic Hadronic  and Nuclear Collisions & \done & \href{https://arxiv.org/pdf/1801.03477}{Link} \\
      \hline
       Progress and Challenges in Small Systems & \done & \href{https://arxiv.org/pdf/2401.09208}{Link} \\
      \hline
    \end{tabular}
  \end{table}
\end{frame}

\begin{frame}
  \frametitle{Recent studies}

  \begin{table}
    \begin{tabular}{|p{0.5\textwidth}|c|p{0.25\textwidth}|}
      \hline
      \textbf{Toolbox} & \textbf{Status} & \textbf{Link} \\
      \hline
       2. Gaussian Processes & \inprogress & \href{https://www.ifsc.usp.br/~roland/pdf/PIQFT_00.pdf}{Lecture Notes}  \\
      \hline
    \end{tabular}
  \end{table}

  \begin{table}
    \begin{tabular}{|p{0.5\textwidth}|c|p{0.25\textwidth}|}
      \hline
      \textbf{Course} & \textbf{Status} & \textbf{Link} \\
      \hline
       Statistics For Applications & \inprogress & \href{https://ocw.mit.edu/courses/18-650-statistics-for-applications-fall-2016/}{MIT OCW 18.650}\\
      \hline
    \end{tabular}
  \end{table}
\end{frame}




% \begin{frame}{Uma imagem}
%    \begin{figure}[h]
%        \centering
%        \includegraphics[width=0.7\textwidth]{img/logo_IFUSP.png}
%        \caption{Legenda da imagem}
%        \label{fig:label_da_imagem}
%    \end{figure}
% \end{frame}

%----------------------------------------------------------------------------------------

% \begin{frame}{Duas imagens}
%    \begin{figure}
%        \centering
%        \begin{subfigure}[b]{0.45\textwidth}
%            \centering
%            \includegraphics[width=\textwidth]{img/logo_IFUSP.png}
%            \caption{Legenda 1}
%            \label{fig:img1}
%        \end{subfigure}
%        \hfill
%        \begin{subfigure}[b]{0.45\textwidth}
%            \centering
%            \includegraphics[width=\textwidth]{img/logo_IFUSP.png}
%            \caption{Legenda 2}
%            \label{fig:img2}
%        \end{subfigure}
%    \end{figure}
% \end{frame}

%----------------------------------------------------------------------------------------
% \begin{frame}{Equações}
%     Equações de Navier-Stokes Forma expandida (3D):
%     \footnotesize
%         \begin{align*}
%             \rho\left(\frac{\partial u}{\partial t} + u\frac{\partial u}{\partial x} + v\frac{\partial u}{\partial y} + w\frac{\partial u}{\partial z}\right) &= -\frac{\partial p}{\partial x} + \mu\left(\frac{\partial^2 u}{\partial x^2} + \frac{\partial^2 u}{\partial y^2} + \frac{\partial^2 u}{\partial z^2}\right) + f_x \\[0.3cm]
%             \rho\left(\frac{\partial v}{\partial t} + u\frac{\partial v}{\partial x} + v\frac{\partial v}{\partial y} + w\frac{\partial v}{\partial z}\right) &= -\frac{\partial p}{\partial y} + \mu\left(\frac{\partial^2 v}{\partial x^2} + \frac{\partial^2 v}{\partial y^2} + \frac{\partial^2 v}{\partial z^2}\right) + f_y \\[0.3cm]
%             \rho\left(\frac{\partial w}{\partial t} + u\frac{\partial w}{\partial x} + v\frac{\partial w}{\partial y} + w\frac{\partial w}{\partial z}\right) &= -\frac{\partial p}{\partial z} + \mu\left(\frac{\partial^2 w}{\partial x^2} + \frac{\partial^2 w}{\partial y^2} + \frac{\partial^2 w}{\partial z^2}\right) + f_z
%         \end{align*}
        
%     onde $\mathbf{v} = (u,v,w)$ é o campo de velocidade, $p$ é a pressão, $\rho$ é a densidade, $\mu$ é a viscosidade dinâmica e $\mathbf{f}$ representa forças externas.
% \end{frame}
\include{sections/section02}
\section{Selected Paper} % Seções são adicionadas para organizar sua apresentação em blocos discretos, todas as seções e subseções são automaticamente exibidas no índice como uma visão geral da apresentação, mas NÃO são exibidas como slides separados.

\begin{frame}{Selected Paper}
\frametitle{Selected Paper}
   \begin{figure}[h]
       \centering
       \includegraphics[width=0.7\textwidth]{img/selected_article.png}
       \caption{The promises and pitfalls of deep kernel learning” (Ober et al., 2021), providing a critical analysis of the strengths, limitations, and failure modes of deep kernel learning models.}
       \label{fig:github_topic_history}
   \end{figure}
\end{frame}

\begin{frame}{References}
    \nocite{*}
    \printbibliography
\end{frame}

% Slide final
% \begin{frame}
%     \begin{center}
%         {\Large Fim da apresentação!}
%     \end{center}
% \end{frame}

\end{document}


